\subsection*{الف}
برای بررسی متعامد بودن، ضرب داخلی‌ها را حساب می‌کنیم:
\begin{flalign*}
s_0=\begin{bmatrix}1\\2\\1\end{bmatrix},\quad
s_1=\begin{bmatrix}1\\0\\-1\end{bmatrix},\quad
s_2=\begin{bmatrix}1\\-1\\1\end{bmatrix} &&
\end{flalign*}
\begin{flalign*}
s_0^{T}s_1 &= 1\cdot 1+2\cdot 0+1\cdot (-1)=0,\\
s_0^{T}s_2 &= 1\cdot 1+2\cdot (-1)+1\cdot 1=0,\\
s_1^{T}s_2 &= 1\cdot 1+0\cdot (-1)+(-1)\cdot 1=0. &&
\end{flalign*}
پس $s_0,s_1,s_2$ متعامد هستند.

\subsection*{ب}
برای یکه‌سازی، اندازه‌ها را می‌یابیم:
\begin{flalign*}
\|s_0\|=\sqrt{1^2+2^2+1^2}=\sqrt{6},\qquad
\|s_1\|=\sqrt{1^2+0^2+(-1)^2}=\sqrt{2},\qquad
\|s_2\|=\sqrt{1^2+(-1)^2+1^2}=\sqrt{3}. &&
\end{flalign*}
پس بردارهای یکه:
\begin{flalign*}
u_0=\frac{s_0}{\sqrt{6}}=\frac{1}{\sqrt6}\begin{bmatrix}1\\2\\1\end{bmatrix},\quad
u_1=\frac{s_1}{\sqrt{2}}=\frac{1}{\sqrt2}\begin{bmatrix}1\\0\\-1\end{bmatrix},\quad
u_2=\frac{s_2}{\sqrt{3}}=\frac{1}{\sqrt3}\begin{bmatrix}1\\-1\\1\end{bmatrix}. &&
\end{flalign*}
ماتریس بردارهای متعامدِ یکه (ستون‌ها پایه‌ها هستند):
\begin{flalign*}
Q=\begin{bmatrix}
\frac{1}{\sqrt6} & \frac{1}{\sqrt2} & \frac{1}{\sqrt3}\\
\frac{2}{\sqrt6} & 0               & -\frac{1}{\sqrt3}\\
\frac{1}{\sqrt6} & -\frac{1}{\sqrt2}& \frac{1}{\sqrt3}
\end{bmatrix}. &&
\end{flalign*}

\subsection*{ج}
چون ستون‌های $Q$ متعامدِ یکه‌اند، داریم:
\begin{flalign*}
Q^{T}Q=I \quad \Rightarrow \quad Q^{-1}=Q^{T}. &&
\end{flalign*}
پس $Q$ یک ماتریس تبدیلِ \lr{orthogonal} برای $s_0,s_1,s_2$ است (و تبدیل مختصات معمولاً با $Q^{T}$ انجام می‌شود).

\subsection*{د}
برای $f=\begin{bmatrix}3\\-6\\5\end{bmatrix}$، تبدیل (مختصات در پایه‌ی متعامدِ یکه) برابر است با:
\begin{flalign*}
\widehat{f}=Q^{T}f=
\begin{bmatrix}
u_0^{T}f\\
u_1^{T}f\\
u_2^{T}f
\end{bmatrix}. &&
\end{flalign*}
\begin{flalign*}
u_0^{T}f &= \frac{1}{\sqrt6}(3)+\frac{2}{\sqrt6}(-6)+\frac{1}{\sqrt6}(5)
= \frac{-4}{\sqrt6} = -\frac{2\sqrt6}{3},\\
u_1^{T}f &= \frac{1}{\sqrt2}(3)+0\cdot(-6)-\frac{1}{\sqrt2}(5)
= \frac{-2}{\sqrt2}=-\sqrt2,\\
u_2^{T}f &= \frac{1}{\sqrt3}(3)-\frac{1}{\sqrt3}(-6)+\frac{1}{\sqrt3}(5)
= \frac{14}{\sqrt3}=\frac{14\sqrt3}{3}. &&
\end{flalign*}
بنابراین:
\begin{flalign*}
\widehat{f}=
\begin{bmatrix}
-\frac{4}{\sqrt6}\\[2pt]
-\sqrt2\\[2pt]
\frac{14}{\sqrt3}
\end{bmatrix}. &&
\end{flalign*}


\subsection*{بررسی}
محاسبات پاسخ قبلی را دوباره چک می‌کنیم:
\begin{flalign*}
s_0=\begin{bmatrix}1\\2\\1\end{bmatrix},\quad
s_1=\begin{bmatrix}1\\0\\-1\end{bmatrix},\quad
s_2=\begin{bmatrix}1\\-1\\1\end{bmatrix} &&
\end{flalign*}

\begin{flalign*}
s_0^{T}s_1&=1\cdot 1+2\cdot 0+1\cdot(-1)=0,\\
s_0^{T}s_2&=1\cdot 1+2\cdot(-1)+1\cdot 1=0,\\
s_1^{T}s_2&=1\cdot 1+0\cdot(-1)+(-1)\cdot 1=0
\quad\Rightarrow\quad s_0,s_1,s_2 \ \text{متعامد}. &&
\end{flalign*}

\begin{flalign*}
\|s_0\|=\sqrt6,\qquad \|s_1\|=\sqrt2,\qquad \|s_2\|=\sqrt3 \\
u_0=\frac{1}{\sqrt6}\begin{bmatrix}1\\2\\1\end{bmatrix},\quad
u_1=\frac{1}{\sqrt2}\begin{bmatrix}1\\0\\-1\end{bmatrix},\quad
u_2=\frac{1}{\sqrt3}\begin{bmatrix}1\\-1\\1\end{bmatrix}. &&
\end{flalign*}

اگر ستون‌ها را پایه بگیریم:
\begin{flalign*}
Q=[u_0\ u_1\ u_2]
=
\begin{bmatrix}
\frac{1}{\sqrt6} & \frac{1}{\sqrt2} & \frac{1}{\sqrt3}\\
\frac{2}{\sqrt6} & 0               & -\frac{1}{\sqrt3}\\
\frac{1}{\sqrt6} & -\frac{1}{\sqrt2}& \frac{1}{\sqrt3}
\end{bmatrix},
\qquad Q^{T}Q=I \Rightarrow Q^{-1}=Q^{T}. &&
\end{flalign*}

نکتهٔ قراردادی: ماتریس تبدیلِ مستقیم برای گرفتن ضرایب معمولاً
\lr{(transform coefficients)}
\begin{flalign*}
T=Q^{T}=
\begin{bmatrix}
u_0^{T}\\
u_1^{T}\\
u_2^{T}
\end{bmatrix}
\quad\Rightarrow\quad \widehat{f}=Tf=Q^{T}f. &&
\end{flalign*}

برای $f=\begin{bmatrix}3\\-6\\5\end{bmatrix}$:
\begin{flalign*}
\widehat{f}=
Q^{T}f=
\begin{bmatrix}
u_0^{T}f\\
u_1^{T}f\\
u_2^{T}f
\end{bmatrix}
=
\begin{bmatrix}
\frac{1}{\sqrt6}(3-12+5)\\
\frac{1}{\sqrt2}(3-5)\\
\frac{1}{\sqrt3}(3+6+5)
\end{bmatrix}
=
\begin{bmatrix}
-\frac{4}{\sqrt6}\\
-\sqrt2\\
\frac{14}{\sqrt3}
\end{bmatrix}. &&
\end{flalign*}

(پس از نظر محاسباتی مشکلی وجود ندارد؛ فقط در بخش (ج) ممکن است بسته به تعریف استاد، $Q$ یا $Q^{T}$ را «ماتریس تبدیل» بنامند که هر دو \lr{orthogonal} هستند.)

\subsection*{ب (ادامه)}
ابتدا توجه می‌کنیم که هرچند $s_0,s_1,s_2$ متعامدند، اما \lr{orthonormal} نیستند چون:
\begin{flalign*}
\|s_0\|=\sqrt6\neq 1,\qquad \|s_1\|=\sqrt2\neq 1,\qquad \|s_2\|=\sqrt3\neq 1. &&
\end{flalign*}
پس آن‌ها را یکه می‌کنیم:
\begin{flalign*}
u_0=\frac{1}{\sqrt6}\begin{bmatrix}1\\2\\1\end{bmatrix},\quad
u_1=\frac{1}{\sqrt2}\begin{bmatrix}1\\0\\-1\end{bmatrix},\quad
u_2=\frac{1}{\sqrt3}\begin{bmatrix}1\\-1\\1\end{bmatrix}. &&
\end{flalign*}
اکنون نشان می‌دهیم $\{u_0,u_1,u_2\}$ یک مجموعه‌ی \lr{orthonormal} است:
\begin{flalign*}
u_0^T u_0 &= \frac{1}{6}(1^2+2^2+1^2)=1,\\
u_1^T u_1 &= \frac{1}{2}(1^2+0^2+(-1)^2)=1,\\
u_2^T u_2 &= \frac{1}{3}(1^2+(-1)^2+1^2)=1,\\
u_0^T u_1 &= \frac{1}{\sqrt{12}}(1+0-1)=0,\\
u_0^T u_2 &= \frac{1}{\sqrt{18}}(1-2+1)=0,\\
u_1^T u_2 &= \frac{1}{\sqrt{6}}(1+0-1)=0. &&
\end{flalign*}
پس ماتریس تبدیلِ ستون‌های متعامدِ یکه (ماتریس پایه) برابر است با:
\begin{flalign*}
Q=\big[u_0\ \ u_1\ \ u_2\big]=
\begin{bmatrix}
\frac{1}{\sqrt6} & \frac{1}{\sqrt2} & \frac{1}{\sqrt3}\\
\frac{2}{\sqrt6} & 0               & -\frac{1}{\sqrt3}\\
\frac{1}{\sqrt6} & -\frac{1}{\sqrt2}& \frac{1}{\sqrt3}
\end{bmatrix},\qquad Q^TQ=I. &&
\end{flalign*}

\subsection*{ج}
در بخش (ب) ماتریس $Q=[u_0\;u_1\;u_2]$ را ساختیم که ستون‌های آن متعامدِ یکه هستند؛ پس:
\begin{flalign*}
Q^{T}Q=I \quad \Rightarrow \quad Q^{-1}=Q^{T}. &&
\end{flalign*}
در \lr{orthonormal transform} معمولاً «تبدیل رو به جلو» (گرفتن ضرایب روی پایه‌ها) با قرار دادن بردارهای یکه به صورت سطر انجام می‌شود، یعنی:
\begin{flalign*}
A=\begin{bmatrix}
u_0^{T}\\
u_1^{T}\\
u_2^{T}
\end{bmatrix}=Q^{T}. &&
\end{flalign*}
بنابراین ماتریس تبدیل متعامد برابر است با:
\begin{flalign*}
A
=
\begin{bmatrix}
\frac{1}{\sqrt6} & \frac{2}{\sqrt6} & \frac{1}{\sqrt6}\\
\frac{1}{\sqrt2} & 0               & -\frac{1}{\sqrt2}\\
\frac{1}{\sqrt3} & -\frac{1}{\sqrt3}& \frac{1}{\sqrt3}
\end{bmatrix},\qquad
AA^{T}=I. &&
\end{flalign*}
اگر $f$ بردارِ ورودی باشد، ضرایب تبدیل:
\begin{flalign*}
\widehat{f}=Af, &&
\end{flalign*}
و بازسازی (تبدیل معکوس) چون $A$ متعامد است:
\begin{flalign*}
f=A^{T}\widehat{f}=Q\,\widehat{f}. &&
\end{flalign*}